\documentclass[tikz,border=4pt]{standalone}
% ==========================================================================
% tikz-tcee.tex — Estilos comunes de los gráficos TikZ del temario
% Autor: Víctor Gutiérrez Marcos
%
% Cada gráfico es un fichero standalone en fuentes/<tema>/graficos/:
%
%   \documentclass[tikz,border=4pt]{standalone}
%   % ==========================================================================
% tikz-tcee.tex — Estilos comunes de los gráficos TikZ del temario
% Autor: Víctor Gutiérrez Marcos
%
% Cada gráfico es un fichero standalone en fuentes/<tema>/graficos/:
%
%   \documentclass[tikz,border=4pt]{standalone}
%   % ==========================================================================
% tikz-tcee.tex — Estilos comunes de los gráficos TikZ del temario
% Autor: Víctor Gutiérrez Marcos
%
% Cada gráfico es un fichero standalone en fuentes/<tema>/graficos/:
%
%   \documentclass[tikz,border=4pt]{standalone}
%   \input{../../../latex/tikz-tcee}
%   \begin{document}
%   \begin{tikzpicture}[tcee]
%     \ejes{Q}{P}{6}{5}
%     \draw[curva1] (0.5,4.5) -- (5,0.8) node[etiqueta,right] {$D$};
%     \capa{2}{\draw[curva2] ...;}   % en el vídeo aparece en el paso 2
%   \end{tikzpicture}
%   \end{document}
%
% Paleta: la de la web (morado, dorado, salvia) más tres colores de apoyo
% distinguibles en blanco y negro por el trazo.
% ==========================================================================
\usepackage[T1]{fontenc}
\usepackage{newpxtext,newpxmath}
\usepackage{xcolor}
\usepackage{tikz,pgfplots}
\pgfplotsset{compat=1.18}
\usetikzlibrary{arrows.meta,calc,positioning,patterns,decorations.pathreplacing,intersections,shapes.misc,backgrounds}

\definecolor{tceemorado}{HTML}{5F2987}
\definecolor{tceedorado}{HTML}{B8860B}
\definecolor{tceeverde}{HTML}{3C7D3A}
\definecolor{tceeazul}{HTML}{1F5F99}
\definecolor{tceerojo}{HTML}{B22222}
\definecolor{tceegris}{HTML}{767171}
\definecolor{tceesalvia}{HTML}{E2EFD9}

% Capas para el vídeo: \capa{n}{…} dibuja su contenido solo a partir del paso n.
% En el PDF, la web y Word se ve todo; generar-video.py compila el gráfico una
% vez por paso con \def\tceepaso{k} para que cada elemento aparezca al narrarlo.
\providecommand{\tceepaso}{1000}
\newcommand{\capa}[2]{\ifnum#1>\tceepaso\relax\else#2\fi}

\tikzset{
  tcee/.style={line cap=round,line join=round,font=\small,>={Stealth[length=2.2mm]}},
  eje/.style={->,thick,black},
  curva1/.style={very thick,tceemorado},
  curva2/.style={very thick,tceeazul},
  curva3/.style={very thick,tceeverde},
  curva4/.style={very thick,tceedorado},
  curvanueva/.style={very thick,tceerojo},
  desplazada/.style={very thick,dashed},
  guia/.style={thin,dashed,tceegris},
  punto/.style={circle,fill=black,inner sep=1.3pt},
  etiqueta/.style={font=\small,inner sep=2pt},
  flecha/.style={->,thick,tceerojo},
  area/.style={fill=tceemorado,fill opacity=0.18,draw=none},
  area2/.style={fill=tceedorado,fill opacity=0.22,draw=none},
}

% \ejes{etiqueta x}{etiqueta y}{largo x}{alto y}
\newcommand{\ejes}[4]{%
  \draw[eje] (0,0) -- (#3,0) node[below left=2pt and 0pt] {#1};
  \draw[eje] (0,0) -- (0,#4) node[left] {#2};
  \node[below left] at (0,0) {$0$};}
% \proyeccion{(x,y)}{etq x}{etq y}: líneas guía a los ejes con etiquetas
\newcommand{\proyeccion}[3]{%
  \draw[guia] let \p1=#1 in (\x1,0) node[below] {#2} |- (0,\y1) node[left] {#3};}

%   \begin{document}
%   \begin{tikzpicture}[tcee]
%     \ejes{Q}{P}{6}{5}
%     \draw[curva1] (0.5,4.5) -- (5,0.8) node[etiqueta,right] {$D$};
%     \capa{2}{\draw[curva2] ...;}   % en el vídeo aparece en el paso 2
%   \end{tikzpicture}
%   \end{document}
%
% Paleta: la de la web (morado, dorado, salvia) más tres colores de apoyo
% distinguibles en blanco y negro por el trazo.
% ==========================================================================
\usepackage[T1]{fontenc}
\usepackage{newpxtext,newpxmath}
\usepackage{xcolor}
\usepackage{tikz,pgfplots}
\pgfplotsset{compat=1.18}
\usetikzlibrary{arrows.meta,calc,positioning,patterns,decorations.pathreplacing,intersections,shapes.misc,backgrounds}

\definecolor{tceemorado}{HTML}{5F2987}
\definecolor{tceedorado}{HTML}{B8860B}
\definecolor{tceeverde}{HTML}{3C7D3A}
\definecolor{tceeazul}{HTML}{1F5F99}
\definecolor{tceerojo}{HTML}{B22222}
\definecolor{tceegris}{HTML}{767171}
\definecolor{tceesalvia}{HTML}{E2EFD9}

% Capas para el vídeo: \capa{n}{…} dibuja su contenido solo a partir del paso n.
% En el PDF, la web y Word se ve todo; generar-video.py compila el gráfico una
% vez por paso con \def\tceepaso{k} para que cada elemento aparezca al narrarlo.
\providecommand{\tceepaso}{1000}
\newcommand{\capa}[2]{\ifnum#1>\tceepaso\relax\else#2\fi}

\tikzset{
  tcee/.style={line cap=round,line join=round,font=\small,>={Stealth[length=2.2mm]}},
  eje/.style={->,thick,black},
  curva1/.style={very thick,tceemorado},
  curva2/.style={very thick,tceeazul},
  curva3/.style={very thick,tceeverde},
  curva4/.style={very thick,tceedorado},
  curvanueva/.style={very thick,tceerojo},
  desplazada/.style={very thick,dashed},
  guia/.style={thin,dashed,tceegris},
  punto/.style={circle,fill=black,inner sep=1.3pt},
  etiqueta/.style={font=\small,inner sep=2pt},
  flecha/.style={->,thick,tceerojo},
  area/.style={fill=tceemorado,fill opacity=0.18,draw=none},
  area2/.style={fill=tceedorado,fill opacity=0.22,draw=none},
}

% \ejes{etiqueta x}{etiqueta y}{largo x}{alto y}
\newcommand{\ejes}[4]{%
  \draw[eje] (0,0) -- (#3,0) node[below left=2pt and 0pt] {#1};
  \draw[eje] (0,0) -- (0,#4) node[left] {#2};
  \node[below left] at (0,0) {$0$};}
% \proyeccion{(x,y)}{etq x}{etq y}: líneas guía a los ejes con etiquetas
\newcommand{\proyeccion}[3]{%
  \draw[guia] let \p1=#1 in (\x1,0) node[below] {#2} |- (0,\y1) node[left] {#3};}

%   \begin{document}
%   \begin{tikzpicture}[tcee]
%     \ejes{Q}{P}{6}{5}
%     \draw[curva1] (0.5,4.5) -- (5,0.8) node[etiqueta,right] {$D$};
%     \capa{2}{\draw[curva2] ...;}   % en el vídeo aparece en el paso 2
%   \end{tikzpicture}
%   \end{document}
%
% Paleta: la de la web (morado, dorado, salvia) más tres colores de apoyo
% distinguibles en blanco y negro por el trazo.
% ==========================================================================
\usepackage[T1]{fontenc}
\usepackage{newpxtext,newpxmath}
\usepackage{xcolor}
\usepackage{tikz,pgfplots}
\pgfplotsset{compat=1.18}
\usetikzlibrary{arrows.meta,calc,positioning,patterns,decorations.pathreplacing,intersections,shapes.misc,backgrounds}

\definecolor{tceemorado}{HTML}{5F2987}
\definecolor{tceedorado}{HTML}{B8860B}
\definecolor{tceeverde}{HTML}{3C7D3A}
\definecolor{tceeazul}{HTML}{1F5F99}
\definecolor{tceerojo}{HTML}{B22222}
\definecolor{tceegris}{HTML}{767171}
\definecolor{tceesalvia}{HTML}{E2EFD9}

% Capas para el vídeo: \capa{n}{…} dibuja su contenido solo a partir del paso n.
% En el PDF, la web y Word se ve todo; generar-video.py compila el gráfico una
% vez por paso con \def\tceepaso{k} para que cada elemento aparezca al narrarlo.
\providecommand{\tceepaso}{1000}
\newcommand{\capa}[2]{\ifnum#1>\tceepaso\relax\else#2\fi}

\tikzset{
  tcee/.style={line cap=round,line join=round,font=\small,>={Stealth[length=2.2mm]}},
  eje/.style={->,thick,black},
  curva1/.style={very thick,tceemorado},
  curva2/.style={very thick,tceeazul},
  curva3/.style={very thick,tceeverde},
  curva4/.style={very thick,tceedorado},
  curvanueva/.style={very thick,tceerojo},
  desplazada/.style={very thick,dashed},
  guia/.style={thin,dashed,tceegris},
  punto/.style={circle,fill=black,inner sep=1.3pt},
  etiqueta/.style={font=\small,inner sep=2pt},
  flecha/.style={->,thick,tceerojo},
  area/.style={fill=tceemorado,fill opacity=0.18,draw=none},
  area2/.style={fill=tceedorado,fill opacity=0.22,draw=none},
}

% \ejes{etiqueta x}{etiqueta y}{largo x}{alto y}
\newcommand{\ejes}[4]{%
  \draw[eje] (0,0) -- (#3,0) node[below left=2pt and 0pt] {#1};
  \draw[eje] (0,0) -- (0,#4) node[left] {#2};
  \node[below left] at (0,0) {$0$};}
% \proyeccion{(x,y)}{etq x}{etq y}: líneas guía a los ejes con etiquetas
\newcommand{\proyeccion}[3]{%
  \draw[guia] let \p1=#1 in (\x1,0) node[below] {#2} |- (0,\y1) node[left] {#3};}

\begin{document}
% Panel izquierdo (ADPR):
%   situación 0: x2 = 2,5 - 0,5 x1 (pendiente -p1^0/p2^0), elige x0 = (3,6; 0,7)
%   situación 1: x2 = 4 - 1,5 x1, elige x1 = (2; 1)
%   x1 es asequible en 0 (2,5 - 1 = 1,5 > 1) y x0 no lo es en 1 (4 - 5,4 < 0,7).
% Panel derecho (signo del efecto sustitución):
%   recta inicial x2 = 4 - x1, elige x0 = (2; 2)
%   baja p1 y se compensa la renta à la Slutsky: x2 = 3 - 0,5 x1, pasa por x0
%   nueva elección x1 = (3,6; 1,2), a la derecha de x0.
\begin{tikzpicture}[tcee]
  % ---------------- Panel izquierdo ----------------
  \capa{1}{\ejes{$x_1$}{$x_2$}{6.2}{4.6}}
  \capa{2}{\draw[curva1] (0,2.5) -- (5,0) node[etiqueta,above right=1pt and -2pt] {$p^0$};}
  \capa{3}{\node[punto,label={[etiqueta]above right:$x^0$}] at (3.6,0.7) {};}
  \capa{4}{\node[punto,label={[etiqueta]below left:$x^1$}] at (2,1) {};}
  \capa{5}{\node[etiqueta,align=center,anchor=south] at (3.9,3.3)
      {$x^0$ se revela\\preferida a $x^1$};}
  \capa{6}{\draw[curva2] (0,4) -- (2.667,0) node[etiqueta,above right=1pt and -2pt] {$p^1$};}
  \capa{7}{\node[etiqueta,align=center,anchor=north,font=\footnotesize] at (4.6,2.1)
      {$x^0$ no es asequible\\cuando se elige $x^1$};}
  % ---------------- Panel derecho ----------------
  \begin{scope}[xshift=7.4cm]
    \capa{8}{\ejes{$x_1$}{$x_2$}{6.6}{4.6}}
    \capa{9}{
      \fill[area] (0,0) -- (0,4) -- (4,0) -- cycle;
      \draw[curva1] (0,4) -- (4,0) node[etiqueta,above right=1pt and -2pt] {$p$};
    }
    \capa{10}{
      \node[punto] (X0) at (2,2) {};
      \node[etiqueta,above right] at (X0) {$x^0$};
    }
    % recta compensada que pasa por x0: tramo izquierdo resaltado y tramo derecho
    \capa{11}{
      \draw[curva4,line width=2.4pt] (0,3) -- (2,2);
      \draw[curva3] (2,2) -- (6,0) node[etiqueta,above right=1pt and -2pt] {$p'$};
    }
    \capa{12}{
      \node[etiqueta,align=left,anchor=south west,font=\footnotesize] at (1.75,3.15)
        {ya era asequible\\y no se eligió};
      \draw[thin,tceegris] (1.85,3.2) -- (1.0,2.58);
    }
    \capa{13}{
      \node[punto] (X1) at (3.6,1.2) {};
      \node[etiqueta,above right] at (X1) {$x^1$};
    }
    \capa{14}{
      \draw[guia] (X0) -- (2,0) node[below] {$x_1^0$};
      \draw[guia] (X1) -- (3.6,0) node[below] {$x_1^1$};
      \draw[flecha] (2.1,0.35) -- (3.5,0.35);
    }
    \capa{15}{\node[etiqueta,align=left,anchor=west,font=\footnotesize] at (3.9,2.4)
        {baja $p_1$ y se\\compensa la renta:\\elige a la derecha};}
  \end{scope}
\end{tikzpicture}
\end{document}
